\documentclass[10pt]{article}

\usepackage[utf8]{inputenc}

\usepackage[T1]{fontenc}

\usepackage{amsmath}

\usepackage{amsfonts}

\usepackage{amssymb}

\usepackage[version=4]{mhchem}

\usepackage{stmaryrd}

\usepackage{bbold}

\usepackage{graphicx}

\usepackage[export]{adjustbox}

\graphicspath{ {./images/} }



\begin{document}

В рамках аксиоматич. теории гармонич. пространств Брело все неотрицательные супергармонич. функции являются Э. Ф. для нек-рого стандартного процесса.



См. также Мартина гранича, Ляпунова стохастическая бункция.



Лит.: [1] Х а н т Д ж.-А., Марковские процессы и потенциалы, пер. с англ., М., 1962; [2] ші и р я е в А. Н., Статистический последовательный анализ, 2 изд, , М., 1976, '[3] Д ы нк и н Е. Б., Марковские процессы, М., 1963, [4] G e t oo r R. K., Markov processes: Ray processes and right prosesses,

B., 1975, [5] III y p M. Г , «Матем заметки», 1973, т. 13, № 4, c. $587-96 ;[6]$ M e y e r P.-A, "Ann Inst. Fourier», 1962, t 12,

p. $125-230 ;[7]$ е г о ж е, там же, 1963, t. $13, \mathcal{N} 2$, p. $357-72$.



M.. . IIIyp.



ӘЛЕМЕНТ АНАЛИТИЧЕСКОЙ ФУНКЦИИ - $\mathrm{co}^{-}$ вокупность $(D, f)$ области $D$ на плоскости комплексного переменного $\mathbb{C}$ и аналитич. функции $f(z)$, заданной в $D$ при помощи нек-рого аналитич. аппарата, позволяющего әффективно асуществить аналитич. продолжение $f(z)$ во всю ее область существования как палной аналитической функции. Наиболее простой и чаще всего применяемой формой Ә. а. ф. является к р у г о в о ї э л е м е н т в виде совокупности степенного ряда



\[

f(z)=\sum_{k=0}^{\infty} c_{k}(z-a)^{k}

\]



и его круга сходимости $D=\{z \in \mathbb{C}:|z-a|<R\}$ с дентром $a$ (ц е н т р ә л е м е н т a) і радиусом сходимости $R>0$. Аналитич. продолжение здесь достигается (быть может, многократным) переразложением ряда (1) для различных центров $b,|b-a| \leqslant R$, по формулам вида



\[

\begin{gathered}

f(z)=\sum_{n=0}^{\infty} d_{n}(z-b)^{n}=c_{0}+ \\

+\left[c_{1}(b-a)+c_{1}(z-b)\right]+\left[c_{2}(b-a)^{2}+\right. \\

\left.+2 c_{2}(b-a)(z-b)+c_{2}(z-b)^{2}\right]+\ldots .

\end{gathered}

\]



Јюбой из элементов $(D, f)$ полной аналитич. функции определяет ее однозначно и может быть представлен шосредством круговых элементов $\mathbf{c}$ центрами $a \in D$. В случае бесконечно удаленного центра $a=\infty$ круговой элемент принимает вид



\[

f(z)=\sum_{k=0}^{\infty} c_{k} z^{-k}

\]



с областью сходимости $D=\{z \in \mathbb{C}:|z|>R\}$.



В процессе аналитич. продолжения функция $f(z)$ может оказаться многозначной и могут появиться соответствующие алгебраическим точкаль ветвления т. н. развет в ле нвы е э лементы вида



\[

\begin{gathered}

f(z)=\sum_{k=m}^{\infty} c_{k}(z-a)^{k / v}, \\

f(z)=\sum_{k=m}^{\infty} c_{k} z^{-k / v},

\end{gathered}

\]



где $v>1$, число $v-1$ наз. по о я д в ле в н ости. Разветвленные әлементы обобщают понятие Э. а. ф., последние в этой связи наз. также не р а з в ет в ленны м и (при $\nu=1)$ ре г уля рны ми (при $m \geqslant 0$ ) ә ле мен т а ми.



В качестве простейшего элемента $(D, f)$ аналитич. функции $f(z)$ многих комплексных переменных $z=$ $=\left(z_{1}, \ldots, z_{n}\right), n>1$, можно принять совокупность кратного степенного ряда



\[

\begin{gathered}

f(z)=\sum_{|k|=0}^{\infty} c_{k}(z-a)^{k}= \\

=\sum_{k_{1}=0}^{\infty} \cdots \sum_{k_{n}=0}^{\infty} c_{k_{1}} \ldots c_{k_{n}}\left(z_{1}-a_{1}\right)^{k_{1}} \ldots\left(z_{n}-a_{n}\right)^{k_{n}},

\end{gathered}

\]



где $a=\left(a_{1}, \ldots, a_{n}\right)$ - центр, $\quad|k|=k_{1}+\ldots+k_{n}, \quad c_{k}=$ $=c_{k_{1}} \ldots c_{k_{n}},(z-a)^{k} \cong\left(z_{1}-a_{1}\right)^{k_{1}} \ldots\left(z_{n}-a_{n}\right)^{k_{n}}$, и нек-рого поликруга



\[

D=\left\{z \in \mathbb{C}^{n}:\left|z_{j}-a_{j}\right|<R_{j}, \quad j=1, \ldots, n\right\},

\]



в к-ром ряд (2) абсолютно сходится. Однако, при $n>1$ следует иметь в виду, тто поликруг не является точной областью абсолютной сходимости стешенного ряда.



Понятие Ә. а. ф. весьма близко к понятию ростка аналитич. функции.



Лum.: [1] М а р к у ш е в и ч А. И., Теория аналитических Функций, 2 изд , т $1-2$, М., $1967-68 ;$ [2] В л а д и м и-



\begin{center}

\includegraphics[max width=\textwidth]{2023_07_09_497c9d99346d41d0c420g-1}

\end{center}



ЭЛЕМЕНТАРНАЯ АБЕЛЕВА ГРУППА - абеЛева группа, порядки всех неединичных әлементов к-рой равды одному и тому же простому числу $p$.



ЭЛЕМЕНТ РНА АРИФМЕТИКА О. А Иванова. арифметика форлальная.



ЭЛЕМЕНТАРНАЯ СИСТЕМА АКСИОМ - система аксиом, записанная на языке узкого исчисления предикатов. Системы аксиом арифметики формальной, теории множеств Цермело - Френкеля (см. А ксиоматическая теория множеств), типов теории - примеры Э. с. а. B. Н. Гришин.



ЭЈЕМЕНТАРНАЯ ТЕОРИЯ - совокупность замкнутых формул логики предикатов 1-й ступени. Э. т. Th $(K)$ к л а с с а $K$ алгебраических систем сигнатуры $\Omega$ наз. совокупность всех замкнутых формул логики предикатов 1-й ступени сигнатуры $\Omega$, истинных на всех системах из класса $K$. Если класс $K$ состоит из одной системы $A$, то Э. т. класса $K$ наз. Э. т. с и с т ем ы $A$. Две алгебраич. системы одной сигнатуры наз. э л е м ен т а р о о к ви в а лен тны ми, если их Э. т. совпадают. Алгебраич. система $A$ сигнатуры $\Omega$ наз. м о д е л ь ю Ә. т. $T$ сигнатуры $\Omega$, если все формулы из $T$ истинны в $A$. Э. т. наз. с о в м е с тн о й, өсли она имеет мадели. Совместная Э. т. наз. п о л н о й, если любые две ее модели элементарно әквивалентны. Класс всех моделей Ә. т. $T$ обозначается $\operatorname{Mod}(T)$. Э. т. $T$ наз. р а з $\mathrm{p}$ е ши м о й, если множество формул $\mathrm{Th}(\operatorname{Mod}(T))$ (т. е. совокупность всех логич. следствий из $T$ ) рекурсивно. Класс $K$ алгебраич. систем сигнатуры $\Omega$ наз. а к с и о м а т и 3 и р у ем ы м, если существует такая Э. т. $T$ сигнатуры $\Omega$, что $K=\operatorname{Mod}(T)$. В этом случае $T$ наз. с о в о к у пI н ос т ь ю а к с и о м для $K$. Класс $K$ тогда и только тогда аксиоматизируем, когда $K=\operatorname{Mod}(\mathrm{Th}(K))$. Напр., класс плотно линейно удорядоченных множеств без наименьшего и наибольшего элементов аксиоматизируем, его Э. т. разрешима, любые две системы из этого класса элементарно эквивалентны, значит Э. т. этого класса полна; кроме того, Э. т. рассматриваемого класса конечно аксиоматизируема. Класс конечных циклич. групп не является аксиоматизируемым, однако его Э. т. разрешима и, значит, рекурсивно аксиоматизируема. Имеются примеры конечно аксиоматизируемых неразрешимых Э. т. Такими являются Ә. т. групп, колец, полей и др. Однако полная рекурсивно аксиоматизируемая Э. т. обязательно разрепима. Поәтому для доказательства разрепимости рекурсивно аксиоматизируемой Э. т. достаточно заметить, что эта Э. т. полна.



Известен ряд методов доказательства полноты. Э. т. наз. к а т е г о ри чн о й в мо ш н с т и $\alpha$, если все ее модели мощности $\alpha$ изоморфны. Э. т., категоричная в нек-рой бесконечной мощности и не имеющая конечных моделей, обязательно полна. Напр., Э. т. алгебраически замкнутых полей фиксированной характеристики рекурсивно аксиоматизируема и категорична в каждой несчетной мощности, а конечных моделей не имеет, поэтому эта теория полна и разрешима. В частности, разрешима Э. т. поля комплексных чисел. Две формулы той же сигнатуры, тто и сигнатура теории $T$, ә к в и в а л е н тны в т е о рии $T$, если эти формулы имеют одинаковые свободные переменные и для любой модели $A$ теории $T$ и любого способа ириписывания этим свободным переменным элементов модели $A$ либо обе формулы одновременно истинны при





\end{document}